\documentclass[11pt]{article}
\usepackage[T1]{fontenc}
\usepackage{amsmath}
\usepackage{amssymb}
\usepackage{amsthm}
\usepackage[hidelinks]{hyperref}
\newtheorem{theorem}{Theorem}[section]
\newtheorem{lemma}[theorem]{Lemma}
\newtheorem{corollary}[theorem]{Corollary}
\theoremstyle{definition}
\newtheorem{definition}[theorem]{Definition}
\begin{document}
\begin{center}
{\LARGE Masked observations}
\end{center}
\section{Masked observations}
\label{ll:observations:masked-observations}
\begin{theorem}
\label{ll:observations:white-count-zero}
For every natural number \(t\), \(\operatorname{whitecount}(t, 0) = 0\).
\end{theorem}
\begin{proof}
The goal follows from \(\operatorname{whitecountzero}(t)\).
\end{proof}
\begin{theorem}
\label{ll:observations:white-count-black-step}
For every natural number \(t\) and natural number \(n\), if \(\operatorname{center}(t) = \operatorname{true}\), then \(\operatorname{whitecount}(t, n + 1) = \operatorname{whitecount}(t + 1, n)\).
\end{theorem}
\begin{proof}
The goal follows from \(\operatorname{whitecountblackstep}(t, n)\).
\end{proof}
\begin{theorem}
\label{ll:observations:white-count-white-step}
For every natural number \(t\) and natural number \(n\), if \(\operatorname{center}(t) = \operatorname{false}\), then \(\operatorname{whitecount}(t, n + 1) = \operatorname{whitecount}(t + 1, n) + 1\).
\end{theorem}
\begin{proof}
The goal follows from \(\operatorname{whitecountwhitestep}(t, n)\).
\end{proof}
\begin{theorem}
\label{ll:observations:prefix-agreement-iff-white-trace}
For every natural number \(t1\) and natural number \(t2\) and natural number \(n\), if \(\operatorname{centerwindowagree}(t1, t2, n)\), then \(\operatorname{leftprefixagree}(t1, t2, n)\) if and only if \(\operatorname{whitetraceagree}(t1, t2, n)\).
\end{theorem}
\begin{proof}
The goal follows from \(\operatorname{prefixagreementiffwhitetrace}(t1, t2, n)\).
\end{proof}
\begin{theorem}
\label{ll:observations:compressed-prefix-repeat}
For every natural number \(T\) and natural number \(P\) and natural number \(t\) and natural number \(n\), if \(0 < P\), then if \(\operatorname{centerperiodic}(T, P)\), then if \(T \le t\), then there exists a natural number \(a\) such that there exists a natural number \(b\) such that \(a < b\) and \(b \le \operatorname{twopow}(\operatorname{whitecount}(t, n)) \land \operatorname{leftprefixagree}(t + a \cdot P, t + b \cdot P, n)\).
\end{theorem}
\begin{proof}
The goal follows from \(\operatorname{compressedprefixrepeat}(T, P, t, n)\).
\end{proof}
\begin{theorem}
\label{ll:observations:separated-prefix-bound}
For every natural number \(T\) and natural number \(P\) and natural number \(t\) and natural number \(n\) and natural number \(M\), if \(0 < P\), then if \(\operatorname{centerperiodic}(T, P)\), then if \(T \le t\), then if \(\operatorname{prefixseparated}(t, P, n, M)\), then \(M \le \operatorname{twopow}(\operatorname{whitecount}(t, n))\).
\end{theorem}
\begin{proof}
The goal follows from \(\operatorname{separatedprefixbound}(T, P, t, n, M)\).
\end{proof}
\begin{theorem}
\label{ll:observations:edge-information-bound}
For every natural number \(T\) and natural number \(P\) and natural number \(t\) and natural number \(n\), if \(0 < P\), then if \(\operatorname{centerperiodic}(T, P)\), then if \(T \le t\), then \(n < t + \operatorname{twopow}(\operatorname{whitecount}(t, n)) \cdot P\).
\end{theorem}
\begin{proof}
The goal follows from \(\operatorname{edgeinformationbound}(T, P, t, n)\).
\end{proof}
\begin{theorem}
\label{ll:observations:information-bound-audit}
For every natural number \(n\) and natural number \(t\) and natural number \(w\) and natural number \(P\), if \(n < (w + 1) \cdot P\), then \(n < t + \operatorname{twopow}(w) \cdot P\).
\end{theorem}
\begin{proof}
The goal follows from \(\operatorname{informationboundaudit}(n, t, w, P)\).
\end{proof}
\section{Central collision recurrence}
\label{ll:observations:central-collisions}
\begin{theorem}
\label{ll:observations:unit-defect-forces-near-collision}
For every natural number \(Q\) and natural number \(t\) and natural number \(r\), if \(\operatorname{leftadjacentdefect}(Q, t) = \operatorname{true}\), then if \(\operatorname{center}(t + r) = \operatorname{true}\), then if \(\operatorname{center}(t + r + Q) = \operatorname{true}\), then \(\operatorname{nearcollision}(Q, t, r + 1, 2 \cdot (r + 1))\).
\end{theorem}
\begin{proof}
The goal follows from \(\operatorname{unitdefectforcesnearcollision}(Q, t, r)\).
\end{proof}
\begin{theorem}
\label{ll:observations:periodic-unit-defect-forces-near-collision}
For every natural number \(T\) and natural number \(P\) and natural number \(m\) and natural number \(t\), if \(0 < P\), then if \(0 < m\), then if \(\operatorname{centerperiodic}(T, P)\), then if \(T \le t\), then if \(\operatorname{leftadjacentdefect}(m \cdot P, t) = \operatorname{true}\), then \(\operatorname{nearcollision}(m \cdot P, t, P, 2 \cdot P)\).
\end{theorem}
\begin{proof}
The goal follows from \(\operatorname{periodicunitdefectforcesnearcollision}(T, P, m, t)\).
\end{proof}
\begin{theorem}
\label{ll:observations:periodic-central-collision-window}
For every natural number \(T\) and natural number \(P\) and natural number \(m\) and natural number \(N\), if \(0 < P\), then if \(0 < m\), then if \(\operatorname{centerperiodic}(T, P)\), then if \(T \le N\), then there exists a natural number \(t\) such that \(N \le t\) and \(t + 1 \le 2 \cdot N + m \cdot P + P \land \operatorname{centralcollision}(m \cdot P, t, 2 \cdot P)\).
\end{theorem}
\begin{proof}
The goal follows from \(\operatorname{periodiccentralcollisionwindow}(T, P, m, N)\).
\end{proof}
\begin{theorem}
\label{ll:observations:periodic-central-collisions}
For every natural number \(T\) and natural number \(P\) and natural number \(m\) and natural number \(N\), if \(0 < P\), then if \(0 < m\), then if \(\operatorname{centerperiodic}(T, P)\), then there exists a natural number \(t\) such that \(N \le t\) and \(\operatorname{centralcollision}(m \cdot P, t, 2 \cdot P)\).
\end{theorem}
\begin{proof}
The goal follows from \(\operatorname{periodiccentralcollisions}(T, P, m, N)\).
\end{proof}
\begin{theorem}
\label{ll:observations:periodic-recurrent-central-column}
For every natural number \(T\) and natural number \(P\) and natural number \(m\), if \(0 < P\), then if \(0 < m\), then if \(\operatorname{centerperiodic}(T, P)\), then \(\operatorname{recurrentcentralcollision}(m \cdot P, 2 \cdot P)\).
\end{theorem}
\begin{proof}
The goal follows from \(\operatorname{periodicrecurrentcentralcolumn}(T, P, m)\).
\end{proof}
\begin{theorem}
\label{ll:observations:black-window-forces-left-collision}
For every natural number \(Q\) and natural number \(b\) and natural number \(r\), if \(0 < r\), then if \(\operatorname{centerwindowagree}(b, b + Q, r + 1)\), then if \(\operatorname{center}(b) = \operatorname{true}\), then if \(\operatorname{leftadjacentdefect}(Q, b + r) = \operatorname{true}\), then there exists a natural number \(s\) such that \(2 \cdot s \le r + 1\) and \(s \le r \land \operatorname{leftcollision}(Q, b + s, r + 1)\).
\end{theorem}
\begin{proof}
The goal follows from \(\operatorname{blackwindowforcesleftcollision}(Q, b, r)\).
\end{proof}
\begin{theorem}
\label{ll:observations:periodic-left-collision-window}
For every natural number \(T\) and natural number \(P\) and natural number \(m\) and natural number \(N\), if \(0 < P\), then if \(0 < m\), then if \(\operatorname{centerperiodic}(T, P)\), then if \(T \le N\), then there exists a natural number \(t\) such that \(N \le t\) and \(t + 1 \le 2 \cdot N + m \cdot P + 2 \cdot P \land \operatorname{leftcollision}(m \cdot P, t, P)\).
\end{theorem}
\begin{proof}
The goal follows from \(\operatorname{periodicleftcollisionwindow}(T, P, m, N)\).
\end{proof}
\begin{theorem}
\label{ll:observations:periodic-left-collisions}
For every natural number \(T\) and natural number \(P\) and natural number \(m\) and natural number \(N\), if \(0 < P\), then if \(0 < m\), then if \(\operatorname{centerperiodic}(T, P)\), then there exists a natural number \(t\) such that \(N \le t\) and \(\operatorname{leftcollision}(m \cdot P, t, P)\).
\end{theorem}
\begin{proof}
The goal follows from \(\operatorname{periodicleftcollisions}(T, P, m, N)\).
\end{proof}
\begin{theorem}
\label{ll:observations:periodic-recurrent-left-column}
For every natural number \(T\) and natural number \(P\) and natural number \(m\), if \(0 < P\), then if \(0 < m\), then if \(\operatorname{centerperiodic}(T, P)\), then \(\operatorname{recurrentleftcollision}(m \cdot P, P)\).
\end{theorem}
\begin{proof}
The goal follows from \(\operatorname{periodicrecurrentleftcolumn}(T, P, m)\).
\end{proof}
\begin{theorem}
\label{ll:observations:collision-code-nonzero}
For every natural number \(Q\) and natural number \(t\) and natural number \(k\), \(\operatorname{collisioncode}(Q, t, k) \ne 0\) if and only if \(\operatorname{collisionbit}(Q, t, k) = \operatorname{true}\).
\end{theorem}
\begin{proof}
The goal follows from \(\operatorname{collisioncodenonzero}(Q, t, k)\).
\end{proof}
\begin{theorem}
\label{ll:observations:collision-code-bound}
For every natural number \(Q\) and natural number \(t\) and natural number \(k\), \(\operatorname{collisioncode}(Q, t, k) \le 6\).
\end{theorem}
\begin{proof}
The goal follows from \(\operatorname{collisioncodebound}(Q, t, k)\).
\end{proof}
\begin{theorem}
\label{ll:observations:collision-code-colours}
For every natural number \(Q\) and natural number \(t\) and natural number \(k\), if \(0 < k\), then \(\operatorname{collisioncode}(Q, t, k) = 1 \leftrightarrow \operatorname{leftcolour}(Q, t, k) = 1 \land \operatorname{leftcolour}(Q, t, k - 1) = 2\) and \((\operatorname{collisioncode}(Q, t, k) = 2 \leftrightarrow \operatorname{leftcolour}(Q, t, k) = 2 \land \operatorname{leftcolour}(Q, t, k - 1) = 1) \land (\operatorname{collisioncode}(Q, t, k) = 3 \leftrightarrow \operatorname{leftcolour}(Q, t, k) = 1 \land \operatorname{leftcolour}(Q, t, k - 1) = 3) \land (\operatorname{collisioncode}(Q, t, k) = 4 \leftrightarrow \operatorname{leftcolour}(Q, t, k) = 3 \land \operatorname{leftcolour}(Q, t, k - 1) = 1) \land (\operatorname{collisioncode}(Q, t, k) = 5 \leftrightarrow \operatorname{leftcolour}(Q, t, k) = 2 \land \operatorname{leftcolour}(Q, t, k - 1) = 3) \land (\operatorname{collisioncode}(Q, t, k) = 6 \leftrightarrow \operatorname{leftcolour}(Q, t, k) = 3 \land \operatorname{leftcolour}(Q, t, k - 1) = 2)\).
\end{theorem}
\begin{proof}
The goal follows from \(\operatorname{collisioncodecolours}(Q, t, k)\).
\end{proof}
\begin{theorem}
\label{ll:observations:advancing-collision-codes}
For every natural number \(Q\) and natural number \(t\) and natural number \(k\), \(\operatorname{advancingatdepth}(Q, t, k)\) if and only if \(\operatorname{collisioncode}(Q, t, k) = 3\) or \(\operatorname{collisioncode}(Q, t, k) = 5\).
\end{theorem}
\begin{proof}
The goal follows from \(\operatorname{advancingcollisioncodes}(Q, t, k)\).
\end{proof}
\begin{theorem}
\label{ll:observations:collision-transmits-left}
For every natural number \(Q\) and natural number \(t\) and natural number \(k\), if \(\operatorname{collisioncode}(Q, t, k) \ne 0\), then \(\operatorname{leftdefect}(Q, t + 1, k) = \operatorname{leftdefect}(Q, t, k + 1)\).
\end{theorem}
\begin{proof}
The goal follows from \(\operatorname{collisiontransmitsleft}(Q, t, k)\).
\end{proof}
\begin{theorem}
\label{ll:observations:collision-gates-agree}
For every natural number \(Q\) and natural number \(t\) and natural number \(k\), if \(\operatorname{collisioncode}(Q, t, k) \ne 0\), then \(\operatorname{leftgatedifference}(Q, t, k) = \operatorname{false}\).
\end{theorem}
\begin{proof}
The goal follows from \(\operatorname{collisiongatesagree}(Q, t, k)\).
\end{proof}
\begin{theorem}
\label{ll:observations:collision-right-transport-table}
For every natural number \(Q\) and natural number \(t\) and natural number \(k\), if \(2 \le k\), then \(\operatorname{collisioncode}(Q, t, k) = 1 \to \operatorname{leftdefect}(Q, t + 1, k - 1) = \operatorname{leftcell}(t, k - 2)\) and \((\operatorname{collisioncode}(Q, t, k) = 2 \to \operatorname{leftdefect}(Q, t + 1, k - 1) = \operatorname{leftcell}(t + Q, k - 2)) \land (\operatorname{collisioncode}(Q, t, k) = 3 \to \operatorname{leftdefect}(Q, t + 1, k - 1) = \operatorname{true}) \land (\operatorname{collisioncode}(Q, t, k) = 4 \to \operatorname{leftdefect}(Q, t + 1, k - 1) = \operatorname{not}(\operatorname{leftcell}(t + Q, k - 2))) \land (\operatorname{collisioncode}(Q, t, k) = 5 \to \operatorname{leftdefect}(Q, t + 1, k - 1) = \operatorname{true}) \land (\operatorname{collisioncode}(Q, t, k) = 6 \to \operatorname{leftdefect}(Q, t + 1, k - 1) = \operatorname{not}(\operatorname{leftcell}(t, k - 2)))\).
\end{theorem}
\begin{proof}
The goal follows from \(\operatorname{collisionrighttransporttable}(Q, t, k)\).
\end{proof}
\begin{theorem}
\label{ll:observations:advancing-right-colour-table}
For every natural number \(Q\) and natural number \(t\) and natural number \(k\), if \(0 < k\), then \(\operatorname{collisioncode}(Q, t, k) = 3 \to \operatorname{leftcolour}(Q, t + 1, k - 1) = 2\) and if \(\operatorname{collisioncode}(Q, t, k) = 5\), then \(\operatorname{leftcolour}(Q, t + 1, k - 1) = 1\).
\end{theorem}
\begin{proof}
The goal follows from \(\operatorname{advancingrightcolourtable}(Q, t, k)\).
\end{proof}
\begin{theorem}
\label{ll:observations:gate-difference-iff-flip}
For every natural number \(Q\) and natural number \(t\) and natural number \(k\), \(\operatorname{leftdefect}(Q, t + 1, k) = \operatorname{not}(\operatorname{leftdefect}(Q, t, k + 1))\) if and only if \(\operatorname{leftgatedifference}(Q, t, k) = \operatorname{true}\).
\end{theorem}
\begin{proof}
The goal follows from \(\operatorname{gatedifferenceiffflip}(Q, t, k)\).
\end{proof}
\begin{theorem}
\label{ll:observations:incoming-defect-erased-iff}
For every natural number \(Q\) and natural number \(t\) and natural number \(k\), if \(\operatorname{leftdefect}(Q, t, k + 1) = \operatorname{true}\), then \(\operatorname{leftdefect}(Q, t + 1, k) = \operatorname{false}\) if and only if \(\operatorname{leftgatedifference}(Q, t, k) = \operatorname{true}\).
\end{theorem}
\begin{proof}
The goal follows from \(\operatorname{incomingdefecterasediff}(Q, t, k)\).
\end{proof}
\begin{theorem}
\label{ll:observations:erasure-has-no-collision}
For every natural number \(Q\) and natural number \(t\) and natural number \(k\), if \(\operatorname{leftdefect}(Q, t, k + 1) = \operatorname{true}\), then if \(\operatorname{leftdefect}(Q, t + 1, k) = \operatorname{false}\), then \(\operatorname{collisioncode}(Q, t, k) = 0\).
\end{theorem}
\begin{proof}
The goal follows from \(\operatorname{erasurehasnocollision}(Q, t, k)\).
\end{proof}
\begin{theorem}
\label{ll:observations:seed-left-colours}
For every natural number \(Q\) and natural number \(k\), if \(0 < k\), then \(\operatorname{leftcolour}(Q, 0, k) = 0\) or \(\operatorname{leftcolour}(Q, 0, k) = 2\).
\end{theorem}
\begin{proof}
The goal follows from \(\operatorname{seedleftcolours}(Q, k)\).
\end{proof}
\begin{theorem}
\label{ll:observations:seed-left-collision-zero}
For every natural number \(Q\) and natural number \(k\), if \(2 \le k\), then \(\operatorname{collisioncode}(Q, 0, k) = 0\).
\end{theorem}
\begin{proof}
The goal follows from \(\operatorname{seedleftcollisionzero}(Q, k)\).
\end{proof}
\begin{theorem}
\label{ll:observations:seed-advancing-collision-birth}
\(\operatorname{collisioncode}(1, 1, 2) = 5\).
\end{theorem}
\begin{proof}
The goal follows from \(\operatorname{seedadvancingcollisionbirth}\).
\end{proof}
\begin{theorem}
\label{ll:observations:seed-gate-erasure}
\(\operatorname{leftdefect}(3, 0, 3) = \operatorname{true}\) and \(\operatorname{leftdefect}(3, 1, 2) = \operatorname{false} \land \operatorname{leftgatedifference}(3, 0, 2) = \operatorname{true}\).
\end{theorem}
\begin{proof}
The goal follows from \(\operatorname{seedgateerasure}\).
\end{proof}
\begin{theorem}
\label{ll:observations:left-defect-black-cone}
For every natural number \(Q\) and natural number \(t\) and natural number \(k\), if \(\operatorname{leftdefect}(Q, t, k) = \operatorname{true}\), then if \(\operatorname{center}(t) = \operatorname{true}\), then if \(\operatorname{center}(t + Q) = \operatorname{true}\), then \(\operatorname{advancingleftcone}(Q, t, k)\).
\end{theorem}
\begin{proof}
The goal follows from \(\operatorname{leftdefectblackcone}(Q, t, k)\).
\end{proof}
\begin{theorem}
\label{ll:observations:black-window-forces-advancing-collision}
For every natural number \(Q\) and natural number \(b\) and natural number \(r\), if \(0 < r\), then if \(\operatorname{centerwindowagree}(b, b + Q, r + 1)\), then if \(\operatorname{center}(b) = \operatorname{true}\), then if \(\operatorname{leftadjacentdefect}(Q, b + r) = \operatorname{true}\), then there exists a natural number \(s\) such that \(2 \cdot s \le r + 1\) and \(s \le r \land \operatorname{advancingleftcollision}(Q, b + s, r + 1)\).
\end{theorem}
\begin{proof}
The goal follows from \(\operatorname{blackwindowforcesadvancingcollision}(Q, b, r)\).
\end{proof}
\begin{theorem}
\label{ll:observations:periodic-advancing-left-collision-window}
For every natural number \(T\) and natural number \(P\) and natural number \(m\) and natural number \(N\), if \(0 < P\), then if \(0 < m\), then if \(\operatorname{centerperiodic}(T, P)\), then if \(T \le N\), then there exists a natural number \(t\) such that \(N \le t\) and \(t + 1 \le 2 \cdot N + m \cdot P + 2 \cdot P \land \operatorname{advancingleftcollision}(m \cdot P, t, P)\).
\end{theorem}
\begin{proof}
The goal follows from \(\operatorname{periodicadvancingleftcollisionwindow}(T, P, m, N)\).
\end{proof}
\begin{theorem}
\label{ll:observations:periodic-recurrent-advancing-left-column}
For every natural number \(T\) and natural number \(P\) and natural number \(m\), if \(0 < P\), then if \(0 < m\), then if \(\operatorname{centerperiodic}(T, P)\), then \(\operatorname{recurrentadvancingleftcollision}(m \cdot P, P)\).
\end{theorem}
\begin{proof}
The goal follows from \(\operatorname{periodicrecurrentadvancingleftcolumn}(T, P, m)\).
\end{proof}
\begin{theorem}
\label{ll:observations:left-defect-false-colours}
For every natural number \(Q\) and natural number \(t\) and natural number \(k\), \(\operatorname{leftdefect}(Q, t, k) = \operatorname{false}\) if and only if \(\operatorname{leftcolour}(Q, t, k) = 0\) or \(\operatorname{leftcolour}(Q, t, k) = 3\).
\end{theorem}
\begin{proof}
The goal follows from \(\operatorname{leftdefectfalsecolours}(Q, t, k)\).
\end{proof}
\begin{theorem}
\label{ll:observations:left-defect-true-colours}
For every natural number \(Q\) and natural number \(t\) and natural number \(k\), \(\operatorname{leftdefect}(Q, t, k) = \operatorname{true}\) if and only if \(\operatorname{leftcolour}(Q, t, k) = 1\) or \(\operatorname{leftcolour}(Q, t, k) = 2\).
\end{theorem}
\begin{proof}
The goal follows from \(\operatorname{leftdefecttruecolours}(Q, t, k)\).
\end{proof}
\begin{theorem}
\label{ll:observations:left-colours-outside-seed-cone}
For every natural number \(Q\) and natural number \(t\) and natural number \(k\), if \(t < k\), then \(\operatorname{leftcolour}(Q, t, k) = 0\) or \(\operatorname{leftcolour}(Q, t, k) = 2\).
\end{theorem}
\begin{proof}
The goal follows from \(\operatorname{leftcoloursoutsideseedcone}(Q, t, k)\).
\end{proof}
\begin{theorem}
\label{ll:observations:collision-code-support}
For every natural number \(Q\) and natural number \(t\) and natural number \(k\), if \(\operatorname{collisioncode}(Q, t, k) \ne 0\), then \(k \le t + 1\).
\end{theorem}
\begin{proof}
The goal follows from \(\operatorname{collisioncodesupport}(Q, t, k)\).
\end{proof}
\begin{theorem}
\label{ll:observations:gate-code-nonzero}
For every natural number \(Q\) and natural number \(t\) and natural number \(k\), \(\operatorname{gatecode}(Q, t, k) \ne 0\) if and only if \(\operatorname{leftgatedifference}(Q, t, k) = \operatorname{true}\).
\end{theorem}
\begin{proof}
The goal follows from \(\operatorname{gatecodenonzero}(Q, t, k)\).
\end{proof}
\begin{theorem}
\label{ll:observations:gate-code-bound}
For every natural number \(Q\) and natural number \(t\) and natural number \(k\), \(\operatorname{gatecode}(Q, t, k) \le 6\).
\end{theorem}
\begin{proof}
The goal follows from \(\operatorname{gatecodebound}(Q, t, k)\).
\end{proof}
\begin{theorem}
\label{ll:observations:gate-code-colours}
For every natural number \(Q\) and natural number \(t\) and natural number \(k\), if \(0 < k\), then \(\operatorname{gatecode}(Q, t, k) = 1 \leftrightarrow \operatorname{leftcolour}(Q, t, k) = 1 \land \operatorname{leftcolour}(Q, t, k - 1) = 0\) and \((\operatorname{gatecode}(Q, t, k) = 2 \leftrightarrow \operatorname{leftcolour}(Q, t, k) = 0 \land \operatorname{leftcolour}(Q, t, k - 1) = 1) \land (\operatorname{gatecode}(Q, t, k) = 3 \leftrightarrow \operatorname{leftcolour}(Q, t, k) = 1 \land \operatorname{leftcolour}(Q, t, k - 1) = 1) \land (\operatorname{gatecode}(Q, t, k) = 4 \leftrightarrow \operatorname{leftcolour}(Q, t, k) = 2 \land \operatorname{leftcolour}(Q, t, k - 1) = 0) \land (\operatorname{gatecode}(Q, t, k) = 5 \leftrightarrow \operatorname{leftcolour}(Q, t, k) = 0 \land \operatorname{leftcolour}(Q, t, k - 1) = 2) \land (\operatorname{gatecode}(Q, t, k) = 6 \leftrightarrow \operatorname{leftcolour}(Q, t, k) = 2 \land \operatorname{leftcolour}(Q, t, k - 1) = 2)\).
\end{theorem}
\begin{proof}
The goal follows from \(\operatorname{gatecodecolours}(Q, t, k)\).
\end{proof}
\begin{theorem}
\label{ll:observations:erasure-event-of-step}
For every natural number \(Q\) and natural number \(t\) and natural number \(k\), if \(\operatorname{leftdefect}(Q, t, k + 1) = \operatorname{true}\), then if \(\operatorname{leftdefect}(Q, t + 1, k) = \operatorname{false}\), then \(\operatorname{erasureevent}(Q, t, k)\).
\end{theorem}
\begin{proof}
The goal follows from \(\operatorname{erasureeventofstep}(Q, t, k)\).
\end{proof}
\begin{theorem}
\label{ll:observations:erasure-black-phase-iff}
For every natural number \(Q\) and natural number \(t\) and natural number \(k\), if \(\operatorname{erasureevent}(Q, t, k)\), then \(\operatorname{leftcolour}(Q, t + 1, k) = 3\) if and only if \(\operatorname{gatecode}(Q, t, k) \le 3\) and \(\operatorname{leftcolour}(Q, t, k + 1) = 2\) or \(4 \le \operatorname{gatecode}(Q, t, k)\) and \(\operatorname{leftcolour}(Q, t, k + 1) = 1\).
\end{theorem}
\begin{proof}
The goal follows from \(\operatorname{erasureblackphaseiff}(Q, t, k)\).
\end{proof}
\begin{theorem}
\label{ll:observations:erasure-outside-seed-cone-white}
For every natural number \(Q\) and natural number \(t\) and natural number \(k\), if \(t + 1 < k\), then if \(\operatorname{erasureevent}(Q, t, k)\), then \(\operatorname{leftcolour}(Q, t + 1, k) = 0\).
\end{theorem}
\begin{proof}
The goal follows from \(\operatorname{erasureoutsideseedconewhite}(Q, t, k)\).
\end{proof}
\begin{theorem}
\label{ll:observations:advancing-source-first-erasure}
For every natural number \(Q\) and natural number \(t\) and natural number \(k\), if \(2 \le k\), then if \(\operatorname{advancingatdepth}(Q, t, k)\), then if \(\operatorname{leftdefect}(Q, t + k, 0) = \operatorname{false}\), then there exists a natural number \(j\) such that \(\operatorname{firstpatherasure}(Q, t, k, j)\).
\end{theorem}
\begin{proof}
The goal follows from \(\operatorname{advancingsourcefirsterasure}(Q, t, k)\).
\end{proof}
\begin{theorem}
\label{ll:observations:periodic-erasure-window}
For every natural number \(T\) and natural number \(P\) and natural number \(m\) and natural number \(N\), if \(0 < P\), then if \(0 < m\), then if \(\operatorname{centerperiodic}(T, P)\), then if \(T \le N\), then there exists a natural number \(u\) such that there exists a natural number \(d\) such that \(N \le u\) and \(u + 1 \le 2 \cdot N + m \cdot P + 3 \cdot P \land d + 2 \le P \land \operatorname{erasureevent}(m \cdot P, u, d)\).
\end{theorem}
\begin{proof}
The goal follows from \(\operatorname{periodicerasurewindow}(T, P, m, N)\).
\end{proof}
\begin{theorem}
\label{ll:observations:periodic-recurrent-source-erasure}
For every natural number \(T\) and natural number \(P\) and natural number \(m\), if \(0 < P\), then if \(0 < m\), then if \(\operatorname{centerperiodic}(T, P)\), then \(\operatorname{recurrentsourceerasure}(m \cdot P, P)\).
\end{theorem}
\begin{proof}
The goal follows from \(\operatorname{periodicrecurrentsourceerasure}(T, P, m)\).
\end{proof}
\begin{theorem}
\label{ll:observations:no-period-iff-no-recurrent-central-collision}
\(\operatorname{noperiod}\) if and only if \(\operatorname{norecurrentcentralcollision}\).
\end{theorem}
\begin{proof}
The goal follows from \(\operatorname{noperiodiffnorecurrentcentralcollision}\).
\end{proof}
\begin{theorem}
\label{ll:observations:no-row-center-period-iff-no-recurrent-central-collision}
\(\operatorname{norowcenterperiod}\) if and only if \(\operatorname{norecurrentcentralcollision}\).
\end{theorem}
\begin{proof}
The goal follows from \(\operatorname{norowcenterperiodiffnorecurrentcentralcollision}\).
\end{proof}
\begin{theorem}
\label{ll:observations:problem-one-iff-no-recurrent-central-collision}
\(\operatorname{problemone}\) if and only if \(\operatorname{norecurrentcentralcollision}\).
\end{theorem}
\begin{proof}
The goal follows from \(\operatorname{problemoneiffnorecurrentcentralcollision}\).
\end{proof}
\begin{theorem}
\label{ll:observations:no-period-of-no-recurrent-central-collision}
If \(\operatorname{norecurrentcentralcollision}\), then \(\operatorname{noperiod}\).
\end{theorem}
\begin{proof}
Assume \(h\).
The goal follows from \(\operatorname{noperiodofnorecurrentcentralcollision}(h)\).
\end{proof}
\begin{theorem}
\label{ll:observations:not-center-periodic-of-no-recurrent-central-collision}
For every natural number \(T\) and natural number \(P\), if \(0 < P\), then if \(\operatorname{norecurrentcentralcollision}\), then not \(\operatorname{centerperiodic}(T, P)\).
\end{theorem}
\begin{proof}
The goal follows from \(\operatorname{notcenterperiodicofnorecurrentcentralcollision}(T, P)\).
\end{proof}
\begin{theorem}
\label{ll:observations:not-eventually-periodic-of-no-recurrent-central-collision}
If \(\operatorname{norecurrentcentralcollision}\), then \(\operatorname{problemone}\).
\end{theorem}
\begin{proof}
Assume \(h\).
The goal follows from \(\operatorname{noteventuallyperiodicofnorecurrentcentralcollision}(h)\).
\end{proof}
\begin{theorem}
\label{ll:observations:pair-phase-step}
For every natural number \(b\) and natural number \(d\) and natural number \(i\), \(\operatorname{pairphase}(\operatorname{wittphi}(b), \operatorname{wittxor}(\operatorname{wittphi}(d), \operatorname{wittshl}(\operatorname{crosscarry}(b, d), 1)), i) = \operatorname{bitvalue}(\operatorname{rowbit}(\operatorname{wittphi}(\operatorname{wittxor}(b, d)), i)) - \operatorname{bitvalue}(\operatorname{rowbit}(\operatorname{wittphi}(b), i))\).
\end{theorem}
\begin{proof}
The goal follows from \(\operatorname{pairphasestep}(b, d, i)\).
\end{proof}
\begin{theorem}
\label{ll:observations:aligned-phase-left}
For every natural number \(Q\) and natural number \(t\) and natural number \(k\), \(\operatorname{alignedphase}(Q, t, t + k + Q) = \operatorname{leftphase}(Q, t, k)\).
\end{theorem}
\begin{proof}
The goal follows from \(\operatorname{alignedphaseleft}(Q, t, k)\).
\end{proof}
\begin{theorem}
\label{ll:observations:phase-zero-iff}
For every natural number \(Q\) and natural number \(t\) and natural number \(k\), \(\operatorname{leftphase}(Q, t, k) = 0\) if and only if \(\operatorname{leftdefect}(Q, t, k) = \operatorname{false}\).
\end{theorem}
\begin{proof}
The goal follows from \(\operatorname{phasezeroiff}(Q, t, k)\).
\end{proof}
\begin{theorem}
\label{ll:observations:phase-unit-iff}
For every natural number \(Q\) and natural number \(t\) and natural number \(k\), \(\operatorname{leftphase}(Q, t, k) = -1\) or \(\operatorname{leftphase}(Q, t, k) = 1\) if and only if \(\operatorname{leftdefect}(Q, t, k) = \operatorname{true}\).
\end{theorem}
\begin{proof}
The goal follows from \(\operatorname{phaseunitiff}(Q, t, k)\).
\end{proof}
\begin{theorem}
\label{ll:observations:left-phase-balance}
For every natural number \(Q\) and natural number \(t\) and natural number \(k\), \(\operatorname{leftphase}(Q, t + 1, k) = \operatorname{leftphase}(Q, t, k + 1) + \operatorname{leftphaseflux}(Q, t, k)\).
\end{theorem}
\begin{proof}
The goal follows from \(\operatorname{leftphasebalance}(Q, t, k)\).
\end{proof}
\begin{theorem}
\label{ll:observations:collision-phase-flip}
For every natural number \(Q\) and natural number \(t\) and natural number \(k\), if \(\operatorname{collisioncode}(Q, t, k) \ne 0\), then \(\operatorname{leftphaseflux}(Q, t, k) = -2 \cdot \operatorname{leftphase}(Q, t, k + 1)\).
\end{theorem}
\begin{proof}
The goal follows from \(\operatorname{collisionphaseflip}(Q, t, k)\).
\end{proof}
\begin{theorem}
\label{ll:observations:advancing-phase-flip}
For every natural number \(Q\) and natural number \(t\) and natural number \(k\), if \(0 < k\), then if \(\operatorname{advancingatdepth}(Q, t, k)\), then \(\operatorname{leftphase}(Q, t + 1, k - 1) = -\operatorname{leftphase}(Q, t, k)\).
\end{theorem}
\begin{proof}
The goal follows from \(\operatorname{advancingphaseflip}(Q, t, k)\).
\end{proof}
\begin{theorem}
\label{ll:observations:gate-phase-cases}
For every natural number \(Q\) and natural number \(t\) and natural number \(k\), \(\operatorname{gatecode}(Q, t, k) = 1 \to \operatorname{leftgatephase}(Q, t, k) = -1\) and \((\operatorname{gatecode}(Q, t, k) = 2 \to \operatorname{leftgatephase}(Q, t, k) = -1) \land (\operatorname{gatecode}(Q, t, k) = 3 \to \operatorname{leftgatephase}(Q, t, k) = -1) \land (\operatorname{gatecode}(Q, t, k) = 4 \to \operatorname{leftgatephase}(Q, t, k) = 1) \land (\operatorname{gatecode}(Q, t, k) = 5 \to \operatorname{leftgatephase}(Q, t, k) = 1) \land (\operatorname{gatecode}(Q, t, k) = 6 \to \operatorname{leftgatephase}(Q, t, k) = 1)\).
\end{theorem}
\begin{proof}
The goal follows from \(\operatorname{gatephasecases}(Q, t, k)\).
\end{proof}
\begin{theorem}
\label{ll:observations:erasure-phase-balance}
For every natural number \(Q\) and natural number \(t\) and natural number \(k\), if \(\operatorname{erasureevent}(Q, t, k)\), then \(\operatorname{leftphase}(Q, t + 1, k) = 0\) and \(\operatorname{leftphaseflux}(Q, t, k) = -\operatorname{leftphase}(Q, t, k + 1) \land (\operatorname{leftphase}(Q, t, k + 1) = -1 \lor \operatorname{leftphase}(Q, t, k + 1) = 1)\).
\end{theorem}
\begin{proof}
The goal follows from \(\operatorname{erasurephasebalance}(Q, t, k)\).
\end{proof}
\begin{theorem}
\label{ll:observations:erasure-reset-product}
For every natural number \(Q\) and natural number \(t\) and natural number \(k\), if \(\operatorname{erasureevent}(Q, t, k)\), then \(\operatorname{leftcolour}(Q, t + 1, k) = 3 \leftrightarrow \operatorname{leftgatephase}(Q, t, k) \cdot \operatorname{leftphase}(Q, t, k + 1) = -1\) and \(\operatorname{leftcolour}(Q, t + 1, k) = 0 \leftrightarrow \operatorname{leftgatephase}(Q, t, k) \cdot \operatorname{leftphase}(Q, t, k + 1) = 1\).
\end{theorem}
\begin{proof}
The goal follows from \(\operatorname{erasureresetproduct}(Q, t, k)\).
\end{proof}
\begin{theorem}
\label{ll:observations:path-phase-balance}
For every natural number \(Q\) and natural number \(t\) and natural number \(k\) and natural number \(n\), \(\operatorname{pathphase}(Q, t, k, n) = \operatorname{leftphase}(Q, t, k) + \operatorname{pathflux}(Q, t, k, n)\).
\end{theorem}
\begin{proof}
The goal follows from \(\operatorname{pathphasebalance}(Q, t, k, n)\).
\end{proof}
\begin{theorem}
\label{ll:observations:seed-left-phase}
For every natural number \(Q\) and natural number \(t\), if \(0 < Q\), then \(\operatorname{leftphase}(Q, t, t + Q) = 1\).
\end{theorem}
\begin{proof}
The goal follows from \(\operatorname{seedleftphase}(Q, t)\).
\end{proof}
\begin{theorem}
\label{ll:observations:seed-center-phase-balance}
For every natural number \(Q\) and natural number \(t\), if \(0 < Q\), then \(\operatorname{leftphase}(Q, t + (t + Q), 0) = 1 + \operatorname{pathflux}(Q, t, t + Q, t + Q)\).
\end{theorem}
\begin{proof}
The goal follows from \(\operatorname{seedcenterphasebalance}(Q, t)\).
\end{proof}
\begin{theorem}
\label{ll:observations:periodic-seed-flux}
For every natural number \(T\) and natural number \(P\) and natural number \(m\) and natural number \(N\), if \(0 < P\), then if \(0 < m\), then if \(\operatorname{centerperiodic}(T, P)\), then if \(T \le N\), then \(\operatorname{pathflux}(m \cdot P, N, N + m \cdot P, N + m \cdot P) = -1\).
\end{theorem}
\begin{proof}
The goal follows from \(\operatorname{periodicseedflux}(T, P, m, N)\).
\end{proof}
\begin{theorem}
\label{ll:observations:seed-phase-cancellation}
\(\operatorname{leftphase}(3, 0, 3) = 1\) and \(\operatorname{leftphase}(3, 1, 2) = 0 \land \operatorname{leftphaseflux}(3, 0, 2) = -1 \land \operatorname{leftgatephase}(3, 0, 2) = 1\).
\end{theorem}
\begin{proof}
The goal follows from \(\operatorname{seedphasecancellation}\).
\end{proof}
\begin{theorem}
\label{ll:observations:rightmost-left-profile}
For every natural number \(Q\) and natural number \(t\) and natural number \(k\) and natural number \(s\), if \(\operatorname{rightmostleft}(Q, t, k)\), then if \(s < k\), then if \(\operatorname{centercleanthrough}(Q, t, s)\), then \(\operatorname{rightmostleft}(Q, t + s, k - s)\).
\end{theorem}
\begin{proof}
The goal follows from \(\operatorname{rightmostleftprofile}(Q, t, k, s)\).
\end{proof}
\begin{theorem}
\label{ll:observations:leading-first-erasure}
For every natural number \(Q\) and natural number \(t\) and natural number \(k\), if \(\operatorname{leadingadvance}(Q, t, k)\), then if \(\operatorname{centercleanthrough}(Q, t, k)\), then \(\operatorname{firstpatherasure}(Q, t, k, k - 2)\) and \(\operatorname{erasureevent}(Q, t + (k - 1), 0) \land (\operatorname{gatecode}(Q, t + (k - 1), 0) = 2 \lor \operatorname{gatecode}(Q, t + (k - 1), 0) = 5) \land \operatorname{center}(t + (k - 1)) = \operatorname{false}\).
\end{theorem}
\begin{proof}
The goal follows from \(\operatorname{leadingfirsterasure}(Q, t, k)\).
\end{proof}
\begin{theorem}
\label{ll:observations:black-window-forces-leading-source}
For every natural number \(Q\) and natural number \(b\) and natural number \(r\), if \(0 < r\), then if \(\operatorname{centerwindowagree}(b, b + Q, r + 1)\), then if \(\operatorname{center}(b) = \operatorname{true}\), then if \(\operatorname{leftadjacentdefect}(Q, b + r) = \operatorname{true}\), then there exists a natural number \(s\) such that there exists a natural number \(k\) such that \(2 \cdot s \le r + 1\) and \(s \le r \land k + s \le r + 1 \land \operatorname{leadingadvance}(Q, b + s, k)\).
\end{theorem}
\begin{proof}
The goal follows from \(\operatorname{blackwindowforcesleadingsource}(Q, b, r)\).
\end{proof}
\begin{theorem}
\label{ll:observations:periodic-leading-source-window}
For every natural number \(T\) and natural number \(P\) and natural number \(m\) and natural number \(N\), if \(0 < P\), then if \(0 < m\), then if \(\operatorname{centerperiodic}(T, P)\), then if \(T \le N\), then there exists a natural number \(t\) such that there exists a natural number \(k\) such that \(N \le t\) and \(t + 1 \le 2 \cdot N + m \cdot P + 2 \cdot P \land k \le P \land \operatorname{leadingadvance}(m \cdot P, t, k)\).
\end{theorem}
\begin{proof}
The goal follows from \(\operatorname{periodicleadingsourcewindow}(T, P, m, N)\).
\end{proof}
\begin{theorem}
\label{ll:observations:periodic-leading-erasure-window}
For every natural number \(T\) and natural number \(P\) and natural number \(m\) and natural number \(N\), if \(0 < P\), then if \(0 < m\), then if \(\operatorname{centerperiodic}(T, P)\), then if \(T \le N\), then there exists a natural number \(t\) such that there exists a natural number \(k\) such that \(N \le t\) and \(t + 1 \le 2 \cdot N + m \cdot P + 2 \cdot P \land k \le P \land \operatorname{leadingadvance}(m \cdot P, t, k) \land \operatorname{firstpatherasure}(m \cdot P, t, k, k - 2) \land (\operatorname{gatecode}(m \cdot P, t + (k - 1), 0) = 2 \lor \operatorname{gatecode}(m \cdot P, t + (k - 1), 0) = 5)\).
\end{theorem}
\begin{proof}
The goal follows from \(\operatorname{periodicleadingerasurewindow}(T, P, m, N)\).
\end{proof}
\begin{theorem}
\label{ll:observations:periodic-recurrent-leading-erasure}
For every natural number \(T\) and natural number \(P\) and natural number \(m\), if \(0 < P\), then if \(0 < m\), then if \(\operatorname{centerperiodic}(T, P)\), then \(\operatorname{recurrentleadingerasure}(m \cdot P, P)\).
\end{theorem}
\begin{proof}
The goal follows from \(\operatorname{periodicrecurrentleadingerasure}(T, P, m)\).
\end{proof}
\begin{theorem}
\label{ll:observations:left-phase-lag-add}
For every natural number \(Q\) and natural number \(R\) and natural number \(t\) and natural number \(k\), \(\operatorname{leftphase}(Q + R, t, k) = \operatorname{leftphase}(Q, t, k) + \operatorname{leftphase}(R, t + Q, k)\).
\end{theorem}
\begin{proof}
The goal follows from \(\operatorname{leftphaselagadd}(Q, R, t, k)\).
\end{proof}
\begin{theorem}
\label{ll:observations:path-flux-lag-add}
For every natural number \(Q\) and natural number \(R\) and natural number \(t\) and natural number \(k\) and natural number \(n\), \(\operatorname{pathflux}(Q + R, t, k, n) = \operatorname{pathflux}(Q, t, k, n) + \operatorname{pathflux}(R, t + Q, k, n)\).
\end{theorem}
\begin{proof}
The goal follows from \(\operatorname{pathfluxlagadd}(Q, R, t, k, n)\).
\end{proof}
\begin{theorem}
\label{ll:observations:leading-phase-reset}
For every natural number \(Q\) and natural number \(t\) and natural number \(k\), if \(\operatorname{leadingadvance}(Q, t, k)\), then if \(\operatorname{centercleanthrough}(Q, t, k)\), then \(\operatorname{leftcolour}(Q, t + k, 0) = 3 \leftrightarrow \operatorname{leftgatephase}(Q, t + (k - 1), 0) \cdot (\operatorname{leftphase}(Q, t, k) + \operatorname{pathflux}(Q, t, k, k - 1)) = -1\) and \(\operatorname{leftcolour}(Q, t + k, 0) = 0 \leftrightarrow \operatorname{leftgatephase}(Q, t + (k - 1), 0) \cdot (\operatorname{leftphase}(Q, t, k) + \operatorname{pathflux}(Q, t, k, k - 1)) = 1\).
\end{theorem}
\begin{proof}
The goal follows from \(\operatorname{leadingphasereset}(Q, t, k)\).
\end{proof}
\begin{theorem}
\label{ll:observations:center-erasure-iff-gate}
For every natural number \(Q\) and natural number \(t\), if \(\operatorname{leftdefect}(Q, t, 0) = \operatorname{false}\), then if \(\operatorname{leftdefect}(Q, t + 1, 0) = \operatorname{false}\), then \(\operatorname{erasureevent}(Q, t, 0)\) if and only if \(\operatorname{gatecode}(Q, t, 0) = 2\) or \(\operatorname{gatecode}(Q, t, 0) = 5\).
\end{theorem}
\begin{proof}
The goal follows from \(\operatorname{centererasureiffgate}(Q, t)\).
\end{proof}
\begin{theorem}
\label{ll:observations:white-block-gate-convex}
For every natural number \(Q\) and natural number \(b\) and natural number \(n\) and natural number \(i\) and natural number \(j\) and natural number \(k\), if \(\operatorname{commonwhiteblock}(Q, b, n)\), then if \(i \le j\), then if \(j \le k\), then if \(k \le n\), then if \(\operatorname{gatecode}(Q, b + i, 0) = 2\) or \(\operatorname{gatecode}(Q, b + i, 0) = 5\), then if \(\operatorname{gatecode}(Q, b + k, 0) = 2\) or \(\operatorname{gatecode}(Q, b + k, 0) = 5\), then \(\operatorname{gatecode}(Q, b + j, 0) = \operatorname{gatecode}(Q, b + i, 0)\).
\end{theorem}
\begin{proof}
The goal follows from \(\operatorname{whiteblockgateconvex}(Q, b, n, i, j, k)\).
\end{proof}
\begin{theorem}
\label{ll:observations:white-block-erasure-convex}
For every natural number \(Q\) and natural number \(b\) and natural number \(n\) and natural number \(i\) and natural number \(j\) and natural number \(k\), if \(\operatorname{commonwhiteblock}(Q, b, n)\), then if \(\operatorname{centercleanthrough}(Q, b, n + 1)\), then if \(i \le j\), then if \(j \le k\), then if \(k \le n\), then if \(\operatorname{erasureevent}(Q, b + i, 0)\), then if \(\operatorname{erasureevent}(Q, b + k, 0)\), then \(\operatorname{erasureevent}(Q, b + j, 0)\) and \(\operatorname{gatecode}(Q, b + j, 0) = \operatorname{gatecode}(Q, b + i, 0)\).
\end{theorem}
\begin{proof}
The goal follows from \(\operatorname{whiteblockerasureconvex}(Q, b, n, i, j, k)\).
\end{proof}
\begin{theorem}
\label{ll:observations:seed-leading-gate-five}
\(\operatorname{leadingadvance}(4, 5, 2)\) and \(\operatorname{centercleanthrough}(4, 5, 2) \land \operatorname{gatecode}(4, 6, 0) = 5\).
\end{theorem}
\begin{proof}
The goal follows from \(\operatorname{seedleadinggatefive}\).
\end{proof}
\begin{theorem}
\label{ll:observations:seed-leading-gate-two}
\(\operatorname{leadingadvance}(7, 7, 4)\) and \(\operatorname{centercleanthrough}(7, 7, 4) \land \operatorname{gatecode}(7, 10, 0) = 2\).
\end{theorem}
\begin{proof}
The goal follows from \(\operatorname{seedleadinggatetwo}\).
\end{proof}
\begin{theorem}
\label{ll:observations:seed-leading-black-reset}
\(\operatorname{leadingadvance}(11, 4, 4)\) and \(\operatorname{centercleanthrough}(11, 4, 4) \land \operatorname{gatecode}(11, 7, 0) = 5 \land \operatorname{leftcolour}(11, 8, 0) = 3\).
\end{theorem}
\begin{proof}
The goal follows from \(\operatorname{seedleadingblackreset}\).
\end{proof}
\begin{theorem}
\label{ll:observations:transfer-recover}
For every truth value \(incomingBit\) and truth value \(laterBit\) and integer \(incoming\) and integer \(gate\), \(\operatorname{recoverphase}(incomingBit, laterBit, \operatorname{transferphase}(incomingBit, laterBit, incoming, gate), gate) = incoming\).
\end{theorem}
\begin{proof}
The goal follows from \(\operatorname{transferrecover}(incomingBit, laterBit, incoming, gate)\).
\end{proof}
\begin{theorem}
\label{ll:observations:transfer-kernel}
For every truth value \(incomingBit\) and truth value \(laterBit\) and integer \(incoming\) and integer \(gate\), \(\operatorname{transferphase}(incomingBit, laterBit, incoming, gate) = 0\) if and only if \(incoming = \operatorname{recoverphase}(incomingBit, laterBit, 0, gate)\).
\end{theorem}
\begin{proof}
The goal follows from \(\operatorname{transferkernel}(incomingBit, laterBit, incoming, gate)\).
\end{proof}
\begin{theorem}
\label{ll:observations:phase-transfer}
For every natural number \(Q\) and natural number \(t\) and natural number \(k\), \(\operatorname{leftphase}(Q, t + 1, k) = \operatorname{transferphase}(\operatorname{leftcell}(t, k + 1), \operatorname{latergate}(Q, t, k), \operatorname{leftphase}(Q, t, k + 1), \operatorname{leftgatephase}(Q, t, k))\).
\end{theorem}
\begin{proof}
The goal follows from \(\operatorname{phasetransfer}(Q, t, k)\).
\end{proof}
\begin{theorem}
\label{ll:observations:phase-recover}
For every natural number \(Q\) and natural number \(t\) and natural number \(k\), \(\operatorname{leftphase}(Q, t, k + 1) = \operatorname{recoverphase}(\operatorname{leftcell}(t, k + 1), \operatorname{latergate}(Q, t, k), \operatorname{leftphase}(Q, t + 1, k), \operatorname{leftgatephase}(Q, t, k))\).
\end{theorem}
\begin{proof}
The goal follows from \(\operatorname{phaserecover}(Q, t, k)\).
\end{proof}
\begin{theorem}
\label{ll:observations:erasure-recovery}
For every natural number \(Q\) and natural number \(t\) and natural number \(k\), if \(\operatorname{erasureevent}(Q, t, k)\), then not \(\operatorname{leftgatephase}(Q, t, k) = 0\) and \(\operatorname{leftphase}(Q, t, k + 1) = \operatorname{recoverphase}(\operatorname{leftcell}(t, k + 1), \operatorname{latergate}(Q, t, k), 0, \operatorname{leftgatephase}(Q, t, k))\).
\end{theorem}
\begin{proof}
The goal follows from \(\operatorname{erasurerecovery}(Q, t, k)\).
\end{proof}
\begin{theorem}
\label{ll:observations:periodic-erasure-orientation}
For every natural number \(T\) and natural number \(P\) and natural number \(m\) and natural number \(N\) and natural number \(g\), if \(0 < P\), then if \(0 < m\), then if \(\operatorname{centerperiodic}(T, P)\), then if \(g = 2\) or \(g = 5\), then there exists a natural number \(t\) such that \(N \le t\) and \(\operatorname{erasureevent}(m \cdot P, t, 0) \land \operatorname{gatecode}(m \cdot P, t, 0) = g\).
\end{theorem}
\begin{proof}
The goal follows from \(\operatorname{periodicerasureorientation}(T, P, m, N, g)\).
\end{proof}
\begin{theorem}
\label{ll:observations:leading-black-before-arrival}
For every natural number \(Q\) and natural number \(t\) and natural number \(k\), if \(\operatorname{leadingadvance}(Q, t, k)\), then if \(\operatorname{centercleanthrough}(Q, t, k)\), then there exists a natural number \(s\) such that \(s + 1 < k\) and \(\operatorname{center}(t + s) = \operatorname{true}\).
\end{theorem}
\begin{proof}
The goal follows from \(\operatorname{leadingblackbeforearrival}(Q, t, k)\).
\end{proof}
\begin{theorem}
\label{ll:observations:black-gap-forces-initial-leading}
For every natural number \(Q\) and natural number \(b\) and natural number \(r\), if \(0 < r\), then if \(\operatorname{centerwindowagree}(b, b + Q, r + 1)\), then if \(\operatorname{center}(b) = \operatorname{true}\), then if \(\operatorname{whitecentergap}(b, r)\), then if \(\operatorname{leftadjacentdefect}(Q, b + r) = \operatorname{true}\), then there exists a natural number \(k\) such that \(k \le r + 1\) and \(\operatorname{leadingadvance}(Q, b, k)\).
\end{theorem}
\begin{proof}
The goal follows from \(\operatorname{blackgapforcesinitialleading}(Q, b, r)\).
\end{proof}
\begin{theorem}
\label{ll:observations:leading-depth-two-black}
For every natural number \(Q\) and natural number \(t\), if \(\operatorname{leadingadvance}(Q, t, 2)\), then if \(\operatorname{centercleanthrough}(Q, t, 2)\), then \(\operatorname{center}(t) = \operatorname{true}\).
\end{theorem}
\begin{proof}
The goal follows from \(\operatorname{leadingdepthtwoblack}(Q, t)\).
\end{proof}
\begin{theorem}
\label{ll:observations:black-gap-erasure-origin}
For every natural number \(Q\) and natural number \(b\) and natural number \(r\) and natural number \(g\), if \(0 < r\), then if \(\operatorname{centerwindowagree}(b, b + Q, r + 1)\), then if \(\operatorname{center}(b) = \operatorname{true}\), then if \(\operatorname{whitecentergap}(b, r)\), then if \(\operatorname{erasureevent}(Q, b + r, 0)\), then if \(\operatorname{gatecode}(Q, b + r, 0) = g\), then there exists a natural number \(k\) such that \(k \le r + 1\) and \(\operatorname{leadingadvance}(Q, b, k) \land \operatorname{firstpatherasure}(Q, b, k, k - 2) \land \operatorname{gatecode}(Q, b + (k - 1), 0) = g\).
\end{theorem}
\begin{proof}
The goal follows from \(\operatorname{blackgaperasureorigin}(Q, b, r, g)\).
\end{proof}
\begin{theorem}
\label{ll:observations:periodic-black-erasure-origin}
For every natural number \(T\) and natural number \(P\) and natural number \(m\) and natural number \(N\) and natural number \(g\), if \(0 < P\), then if \(0 < m\), then if \(\operatorname{centerperiodic}(T, P)\), then if \(g = 2\) or \(g = 5\), then there exists a natural number \(b\) such that there exists a natural number \(k\) such that \(N \le b\) and \(k \le P \land \operatorname{center}(b) = \operatorname{true} \land \operatorname{whitecentergap}(b, k - 1) \land \operatorname{leadingadvance}(m \cdot P, b, k) \land \operatorname{firstpatherasure}(m \cdot P, b, k, k - 2) \land \operatorname{gatecode}(m \cdot P, b + (k - 1), 0) = g\).
\end{theorem}
\begin{proof}
The goal follows from \(\operatorname{periodicblackerasureorigin}(T, P, m, N, g)\).
\end{proof}
\begin{theorem}
\label{ll:observations:right-import-step}
For every natural number \(t\), \(\operatorname{rightcell}(t + 1, 1) = \operatorname{bxor}(\operatorname{center}(t), \operatorname{bxor}(\operatorname{rightcell}(t, 1), \operatorname{rightimport}(t)))\).
\end{theorem}
\begin{proof}
The goal follows from \(\operatorname{rightimportstep}(t)\).
\end{proof}
\begin{theorem}
\label{ll:observations:right-defect-import-balance}
For every natural number \(Q\) and natural number \(t\), if \(\operatorname{center}(t + Q) = \operatorname{center}(t)\), then \(\operatorname{rightdefect}(Q, t + 1, 1) = \operatorname{bxor}(\operatorname{rightdefect}(Q, t, 1), \operatorname{rightimportdifference}(Q, t))\).
\end{theorem}
\begin{proof}
The goal follows from \(\operatorname{rightdefectimportbalance}(Q, t)\).
\end{proof}
\begin{theorem}
\label{ll:observations:black-reset-survival}
For every natural number \(Q\) and natural number \(t\), if \(\operatorname{center}(t) = \operatorname{true}\), then if \(\operatorname{center}(t + Q) = \operatorname{true}\), then if \(\operatorname{rightdefect}(Q, t, 1) = \operatorname{true}\), then if \(\operatorname{rightdefect}(Q, t + 1, 1) = \operatorname{true}\), then \(\operatorname{rightphase}(Q, t + 1, 1) = -\operatorname{rightphase}(Q, t, 1)\).
\end{theorem}
\begin{proof}
The goal follows from \(\operatorname{blackresetsurvival}(Q, t)\).
\end{proof}
\begin{theorem}
\label{ll:observations:black-reset-import}
For every natural number \(Q\) and natural number \(t\), if \(\operatorname{center}(t) = \operatorname{true}\), then if \(\operatorname{center}(t + Q) = \operatorname{true}\), then if \(\operatorname{rightdefect}(Q, t, 1) = \operatorname{false}\), then if \(\operatorname{rightdefect}(Q, t + 1, 1) = \operatorname{true}\), then \(\operatorname{rightcell}(t, 1) = \operatorname{false}\) and \(\operatorname{rightcell}(t + Q, 1) = \operatorname{false} \land \operatorname{rightdefect}(Q, t, 2) = \operatorname{true} \land \operatorname{rightphase}(Q, t + 1, 1) = -\operatorname{rightphase}(Q, t, 2)\).
\end{theorem}
\begin{proof}
The goal follows from \(\operatorname{blackresetimport}(Q, t)\).
\end{proof}
\begin{theorem}
\label{ll:observations:periodic-right-import-orientation}
For every natural number \(T\) and natural number \(P\) and natural number \(m\) and natural number \(N\) and truth value \(before\), if \(0 < P\), then if \(0 < m\), then if \(\operatorname{centerperiodic}(T, P)\), then there exists a natural number \(t\) such that \(N \le t\) and \(\operatorname{rightimport}(t) = before \land \operatorname{rightimport}(t + m \cdot P) = \operatorname{not}(before) \land \operatorname{rightdefect}(m \cdot P, t + 1, 1) = \operatorname{not}(\operatorname{rightdefect}(m \cdot P, t, 1))\).
\end{theorem}
\begin{proof}
The goal follows from \(\operatorname{periodicrightimportorientation}(T, P, m, N, before)\).
\end{proof}
\begin{theorem}
\label{ll:observations:right-head-birth}
For every natural number \(Q\) and natural number \(t\), if \(\operatorname{center}(t + Q) = \operatorname{center}(t)\), then if \(\operatorname{rightdefect}(Q, t, 1) = \operatorname{false}\), then if \(\operatorname{rightdefect}(Q, t + 1, 1) = \operatorname{true}\), then \(\operatorname{rightcell}(t, 1) = \operatorname{false}\) and \(\operatorname{rightcell}(t + Q, 1) = \operatorname{false} \land \operatorname{rightdefect}(Q, t, 2) = \operatorname{true} \land \operatorname{rightimportdifference}(Q, t) = \operatorname{true} \land \operatorname{rightphase}(Q, t + 1, 1) = (1 - 2 \cdot \operatorname{bitvalue}(\operatorname{center}(t))) \cdot \operatorname{rightphase}(Q, t, 2)\).
\end{theorem}
\begin{proof}
The goal follows from \(\operatorname{rightheadbirth}(Q, t)\).
\end{proof}
\begin{theorem}
\label{ll:observations:head-import-path-phase}
For every natural number \(Q\) and natural number \(start\) and natural number \(finish\), if \(\operatorname{headimportpath}(Q, start, finish)\), then if \(\operatorname{centercleanthrough}(Q, start, finish - start)\), then \(\operatorname{rightphase}(Q, finish, 1) = \operatorname{centersign}(start, finish - start) \cdot \operatorname{rightphase}(Q, start, 2)\).
\end{theorem}
\begin{proof}
The goal follows from \(\operatorname{headimportpathphase}(Q, start, finish)\).
\end{proof}
\begin{theorem}
\label{ll:observations:white-center-phase-matching}
For every natural number \(Q\) and natural number \(t\), if \(\operatorname{leftdefect}(Q, t, 0) = \operatorname{false}\), then if \(\operatorname{leftdefect}(Q, t + 1, 0) = \operatorname{false}\), then if \(\operatorname{center}(t) = \operatorname{false}\), then \(\operatorname{leftphase}(Q, t, 1) = (1 - 2 \cdot \operatorname{bitvalue}(\operatorname{center}(t + 1))) \cdot \operatorname{rightphase}(Q, t, 1)\).
\end{theorem}
\begin{proof}
The goal follows from \(\operatorname{whitecenterphasematching}(Q, t)\).
\end{proof}
\begin{theorem}
\label{ll:observations:periodic-coupled-cancellation}
For every natural number \(T\) and natural number \(P\) and natural number \(m\) and natural number \(N\) and natural number \(g\), if \(0 < P\), then if \(0 < m\), then if \(\operatorname{centerperiodic}(T, P)\), then if \(g = 2\) or \(g = 5\), then there exists a natural number \(start\) such that there exists a natural number \(b\) such that there exists a natural number \(k\) such that \(N \le start\) and \(N \le b \land k \le P \land \operatorname{center}(b) = \operatorname{true} \land \operatorname{whitecentergap}(b, k - 1) \land \operatorname{leadingadvance}(m \cdot P, b, k) \land \operatorname{firstpatherasure}(m \cdot P, b, k, k - 2) \land \operatorname{gatecode}(m \cdot P, b + (k - 1), 0) = g \land \operatorname{headimportpath}(m \cdot P, start, b + (k - 1)) \land \operatorname{leftphase}(m \cdot P, b, k) + \operatorname{pathflux}(m \cdot P, b, k, k - 1) = (1 - 2 \cdot \operatorname{bitvalue}(\operatorname{center}(b + k))) \cdot \operatorname{centersign}(start, b + (k - 1) - start) \cdot \operatorname{rightphase}(m \cdot P, start, 2)\).
\end{theorem}
\begin{proof}
The goal follows from \(\operatorname{periodiccoupledcancellation}(T, P, m, N, g)\).
\end{proof}
\begin{theorem}
\label{ll:observations:leading-arrival-phase-eq-gate-phase}
For every natural number \(Q\) and natural number \(b\) and natural number \(k\), if \(\operatorname{leadingadvance}(Q, b, k)\), then if \(\operatorname{centercleanthrough}(Q, b, k)\), then \(\operatorname{leftphase}(Q, b, k) + \operatorname{pathflux}(Q, b, k, k - 1) = (1 - 2 \cdot \operatorname{bitvalue}(\operatorname{center}(b + k))) \cdot \operatorname{leftgatephase}(Q, b + (k - 1), 0)\).
\end{theorem}
\begin{proof}
The goal follows from \(\operatorname{leadingarrivalphaseeqgatephase}(Q, b, k)\).
\end{proof}
\begin{theorem}
\label{ll:observations:head-import-path-gate-phase}
For every natural number \(Q\) and natural number \(start\) and natural number \(finish\), if \(\operatorname{headimportpath}(Q, start, finish)\), then if \(\operatorname{centercleanthrough}(Q, start, finish - start)\), then if \(\operatorname{center}(finish) = \operatorname{false}\), then \(\operatorname{leftgatephase}(Q, finish, 0) = \operatorname{centersign}(start, finish - start) \cdot \operatorname{rightphase}(Q, start, 2)\).
\end{theorem}
\begin{proof}
The goal follows from \(\operatorname{headimportpathgatephase}(Q, start, finish)\).
\end{proof}
\begin{theorem}
\label{ll:observations:head-import-path-rejects-negative-gate}
For every natural number \(Q\) and natural number \(start\) and natural number \(finish\), if \(\operatorname{headimportpath}(Q, start, finish)\), then if \(\operatorname{centercleanthrough}(Q, start, finish - start)\), then if \(\operatorname{center}(finish) = \operatorname{false}\), then not \(\operatorname{leftgatephase}(Q, finish, 0) = -(\operatorname{centersign}(start, finish - start) \cdot \operatorname{rightphase}(Q, start, 2))\).
\end{theorem}
\begin{proof}
The goal follows from \(\operatorname{headimportpathrejectsnegativegate}(Q, start, finish)\).
\end{proof}
\begin{theorem}
\label{ll:observations:periodic-coupled-gate-phase-sign}
For every natural number \(T\) and natural number \(P\) and natural number \(m\) and natural number \(N\) and natural number \(g\), if \(0 < P\), then if \(0 < m\), then if \(\operatorname{centerperiodic}(T, P)\), then if \(g = 2\) or \(g = 5\), then there exists a natural number \(start\) such that there exists a natural number \(b\) such that there exists a natural number \(k\) such that \(N \le start\) and \(N \le b \land k \le P \land \operatorname{center}(b) = \operatorname{true} \land \operatorname{whitecentergap}(b, k - 1) \land \operatorname{leadingadvance}(m \cdot P, b, k) \land \operatorname{firstpatherasure}(m \cdot P, b, k, k - 2) \land \operatorname{gatecode}(m \cdot P, b + (k - 1), 0) = g \land \operatorname{headimportpath}(m \cdot P, start, b + (k - 1)) \land \operatorname{leftgatephase}(m \cdot P, b + (k - 1), 0) = \operatorname{centersign}(start, b + (k - 1) - start) \cdot \operatorname{rightphase}(m \cdot P, start, 2)\).
\end{theorem}
\begin{proof}
The goal follows from \(\operatorname{periodiccoupledgatephasesign}(T, P, m, N, g)\).
\end{proof}
\begin{theorem}
\label{ll:observations:periodic-doubled-seed-flux-split}
For every natural number \(T\) and natural number \(P\) and natural number \(m\) and natural number \(N\), if \(0 < P\), then if \(0 < m\), then if \(\operatorname{centerperiodic}(T, P)\), then if \(T \le N\), then \(\operatorname{pathflux}(m \cdot P + m \cdot P, N, N + (m \cdot P + m \cdot P), N + (m \cdot P + m \cdot P)) = -1\) and \(\operatorname{pathflux}(m \cdot P, N, N + (m \cdot P + m \cdot P), N + (m \cdot P + m \cdot P)) = 0 \land \operatorname{pathflux}(m \cdot P, N + m \cdot P, N + (m \cdot P + m \cdot P), N + (m \cdot P + m \cdot P)) = -1\).
\end{theorem}
\begin{proof}
The goal follows from \(\operatorname{periodicdoubledseedfluxsplit}(T, P, m, N)\).
\end{proof}
\begin{theorem}
\label{ll:observations:seed-coupled-gate-sign}
\(\operatorname{leadingadvance}(11, 4, 4)\) and \(\operatorname{centercleanthrough}(11, 4, 4) \land \operatorname{headimportpath}(11, 6, 7) \land \operatorname{centercleanthrough}(11, 6, 1) \land \operatorname{leftgatephase}(11, 7, 0) = 1 \land \operatorname{centersign}(6, 1) = 1 \land \operatorname{rightphase}(11, 6, 2) = 1\).
\end{theorem}
\begin{proof}
The goal follows from \(\operatorname{seedcoupledgatesign}\).
\end{proof}
\begin{theorem}
\label{ll:observations:right-defect-has-ancestor}
For every natural number \(Q\) and natural number \(t\) and natural number \(j\), if \(\operatorname{rightdefect}(Q, t + 1, j + 1) = \operatorname{true}\), then \(\operatorname{rightdefect}(Q, t, j) = \operatorname{true}\) or \(\operatorname{rightdefect}(Q, t, j + 1) = \operatorname{true} \lor \operatorname{rightdefect}(Q, t, j + 2) = \operatorname{true}\).
\end{theorem}
\begin{proof}
The goal follows from \(\operatorname{rightdefecthasancestor}(Q, t, j)\).
\end{proof}
\begin{theorem}
\label{ll:observations:right-lineage-of-clean-center}
For every natural number \(Q\) and natural number \(t\) and natural number \(j\) and natural number \(T\), if \(1 \le j\), then if \(T \le t\), then if \(\operatorname{rightdefect}(Q, t, j) = \operatorname{true}\), then if \(\operatorname{centercleanthrough}(Q, T, t - T)\), then \(\operatorname{rightlineage}(Q, t, j, T)\).
\end{theorem}
\begin{proof}
The goal follows from \(\operatorname{rightlineageofcleancenter}(Q, t, j, T)\).
\end{proof}
\begin{theorem}
\label{ll:observations:left-lineage-of-clean-center}
For every natural number \(Q\) and natural number \(t\) and natural number \(k\) and natural number \(T\), if \(1 \le k\), then if \(T \le t\), then if \(\operatorname{leftdefect}(Q, t, k) = \operatorname{true}\), then if \(\operatorname{centercleanthrough}(Q, T, t - T)\), then \(\operatorname{leftlineage}(Q, t, k, T)\).
\end{theorem}
\begin{proof}
The goal follows from \(\operatorname{leftlineageofcleancenter}(Q, t, k, T)\).
\end{proof}
\begin{theorem}
\label{ll:observations:right-lineage-at-start}
For every natural number \(Q\) and natural number \(t\) and natural number \(j\) and natural number \(T\), if \(T \le t\), then if \(\operatorname{rightlineage}(Q, t, j, T)\), then there exists a natural number \(i\) such that \(1 \le i\) and \(i \le j + (t - T) \land j \le i + (t - T) \land \operatorname{rightdefect}(Q, T, i) = \operatorname{true}\).
\end{theorem}
\begin{proof}
The goal follows from \(\operatorname{rightlineageatstart}(Q, t, j, T)\).
\end{proof}
\begin{theorem}
\label{ll:observations:left-lineage-at-start}
For every natural number \(Q\) and natural number \(t\) and natural number \(k\) and natural number \(T\), if \(T \le t\), then if \(\operatorname{leftlineage}(Q, t, k, T)\), then there exists a natural number \(i\) such that \(1 \le i\) and \(i \le k + (t - T) \land k \le i + (t - T) \land \operatorname{leftdefect}(Q, T, i) = \operatorname{true}\).
\end{theorem}
\begin{proof}
The goal follows from \(\operatorname{leftlineageatstart}(Q, t, k, T)\).
\end{proof}
\begin{theorem}
\label{ll:observations:right-lineage-seed}
For every natural number \(Q\) and natural number \(t\) and natural number \(j\), if \(\operatorname{rightlineage}(Q, t, j, 0)\), then there exists a natural number \(i\) such that \(1 \le i\) and \(i \le Q \land i \le j + t \land j \le i + t \land \operatorname{rightdefect}(Q, 0, i) = \operatorname{true} \land \operatorname{rightcell}(Q, i) = \operatorname{true} \land \operatorname{rightphase}(Q, 0, i) = 1\).
\end{theorem}
\begin{proof}
The goal follows from \(\operatorname{rightlineageseed}(Q, t, j)\).
\end{proof}
\begin{theorem}
\label{ll:observations:left-lineage-seed}
For every natural number \(Q\) and natural number \(t\) and natural number \(k\), if \(\operatorname{leftlineage}(Q, t, k, 0)\), then there exists a natural number \(i\) such that \(1 \le i\) and \(i \le Q \land i \le k + t \land k \le i + t \land \operatorname{leftdefect}(Q, 0, i) = \operatorname{true} \land \operatorname{leftcell}(Q, i) = \operatorname{true} \land \operatorname{leftphase}(Q, 0, i) = 1\).
\end{theorem}
\begin{proof}
The goal follows from \(\operatorname{leftlineageseed}(Q, t, k)\).
\end{proof}
\begin{theorem}
\label{ll:observations:exists-right-front}
For every natural number \(Q\) and natural number \(t\), if \(0 < Q\), then there exists a natural number \(y\) such that \(y \le t + Q\) and \(\operatorname{rightfront}(Q, t, y)\).
\end{theorem}
\begin{proof}
The goal follows from \(\operatorname{existsrightfront}(Q, t)\).
\end{proof}
\begin{theorem}
\label{ll:observations:right-front-lower}
For every natural number \(Q\) and natural number \(t\) and natural number \(y\) and natural number \(z\), if \(\operatorname{leftdefect}(Q, t, 0) = \operatorname{false}\), then if \(\operatorname{rightfront}(Q, t, y)\), then if \(\operatorname{rightfront}(Q, t + 1, z)\), then \(y \le z + 1\).
\end{theorem}
\begin{proof}
The goal follows from \(\operatorname{rightfrontlower}(Q, t, y, z)\).
\end{proof}
\begin{theorem}
\label{ll:observations:right-front-descent}
For every natural number \(Q\) and natural number \(t\) and natural number \(y\), if \(\operatorname{leftdefect}(Q, t, 0) = \operatorname{false}\), then if \(\operatorname{rightfront}(Q, t, y)\), then if \(2 \le y\), then \(\operatorname{rightdefect}(Q, t + 1, y - 1) = \operatorname{not}(\operatorname{rightcell}(t, y - 1))\).
\end{theorem}
\begin{proof}
The goal follows from \(\operatorname{rightfrontdescent}(Q, t, y)\).
\end{proof}
\begin{theorem}
\label{ll:observations:right-front-lockstep}
For every natural number \(Q\) and natural number \(t\) and natural number \(y\) and natural number \(z\), if \(\operatorname{leftdefect}(Q, t, 0) = \operatorname{false}\), then if \(\operatorname{rightfront}(Q, t, y)\), then if \(2 \le y\), then if \(\operatorname{rightfront}(Q, t + 1, z)\), then if \(z + 1 \le y\), then \(z = y - 1\) and \(\operatorname{rightcell}(t, y - 1) = \operatorname{false}\).
\end{theorem}
\begin{proof}
The goal follows from \(\operatorname{rightfrontlockstep}(Q, t, y, z)\).
\end{proof}
\begin{theorem}
\label{ll:observations:leading-head-black-before-arrival}
For every natural number \(Q\) and natural number \(b\) and natural number \(k\) and natural number \(s\), if \(\operatorname{rightmostleft}(Q, b, k)\), then if \(\operatorname{centercleanthrough}(Q, b, k)\), then if \(s + 1 < k\), then if \(\operatorname{rightdefect}(Q, b + s, 1) = \operatorname{true}\), then \(\operatorname{center}(b + s) = \operatorname{true}\).
\end{theorem}
\begin{proof}
The goal follows from \(\operatorname{leadingheadblackbeforearrival}(Q, b, k, s)\).
\end{proof}
\begin{theorem}
\label{ll:observations:coupled-import-born-at-arrival}
For every natural number \(Q\) and natural number \(b\) and natural number \(k\) and natural number \(start\), if \(3 \le k\), then if \(\operatorname{leadingadvance}(Q, b, k)\), then if \(\operatorname{centercleanthrough}(Q, b, k)\), then if \(\operatorname{whitecentergap}(b, k - 1)\), then if \(\operatorname{headimportpath}(Q, start, b + (k - 1))\), then \(start = b + (k - 2)\).
\end{theorem}
\begin{proof}
The goal follows from \(\operatorname{coupledimportbornatarrival}(Q, b, k, start)\).
\end{proof}
\begin{theorem}
\label{ll:observations:coupled-right-front-mirror-at}
For every natural number \(Q\) and natural number \(b\) and natural number \(k\) and natural number \(s\), if \(0 < Q\), then if \(3 \le k\), then if \(\operatorname{leadingadvance}(Q, b, k)\), then if \(\operatorname{centercleanthrough}(Q, b, k)\), then if \(\operatorname{whitecentergap}(b, k - 1)\), then if \(1 \le s\), then if \(s + 1 < k\), then there exists a natural number \(y\) such that \(\operatorname{rightdefect}(Q, b + s, 1) = \operatorname{false}\) and \(\operatorname{rightfront}(Q, b + s, y) \land 2 \le y \land y \le k - s\).
\end{theorem}
\begin{proof}
The goal follows from \(\operatorname{coupledrightfrontmirrorat}(Q, b, k, s)\).
\end{proof}
\begin{theorem}
\label{ll:observations:coupled-white-staircase-at}
For every natural number \(Q\) and natural number \(b\) and natural number \(k\) and natural number \(s\), if \(0 < Q\), then if \(3 \le k\), then if \(\operatorname{leadingadvance}(Q, b, k)\), then if \(\operatorname{centercleanthrough}(Q, b, k)\), then if \(\operatorname{whitecentergap}(b, k - 1)\), then if \(1 \le s\), then if \(s + 2 < k\), then if \(\operatorname{rightfront}(Q, b + s, k - s)\), then \(\operatorname{rightcell}(b + s, k - s - 1) = \operatorname{false}\) and \(\operatorname{rightfront}(Q, b + (s + 1), k - (s + 1))\).
\end{theorem}
\begin{proof}
The goal follows from \(\operatorname{coupledwhitestaircaseat}(Q, b, k, s)\).
\end{proof}
\begin{theorem}
\label{ll:observations:periodic-coupled-import-ancestry}
For every natural number \(T\) and natural number \(P\) and natural number \(m\) and natural number \(N\) and natural number \(g\), if \(0 < P\), then if \(0 < m\), then if \(\operatorname{centerperiodic}(T, P)\), then if \(g = 2\) or \(g = 5\), then there exists a natural number \(start\) such that there exists a natural number \(b\) such that there exists a natural number \(k\) such that there exists a natural number \(y\) such that \(N \le start\) and \(N \le b \land T \le b \land T \le start \land 2 \le k \land k \le P \land \operatorname{center}(b) = \operatorname{true} \land \operatorname{whitecentergap}(b, k - 1) \land \operatorname{leadingadvance}(m \cdot P, b, k) \land \operatorname{gatecode}(m \cdot P, b + (k - 1), 0) = g \land \operatorname{headimportpath}(m \cdot P, start, b + (k - 1)) \land (3 \le k \to start = b + (k - 2)) \land (3 \le k \to \operatorname{rightfront}(m \cdot P, b + 1, y) \land 2 \le y \land y \le k - 1) \land \operatorname{rightlineage}(m \cdot P, start, 2, T) \land \operatorname{leftlineage}(m \cdot P, b, k, T)\).
\end{theorem}
\begin{proof}
The goal follows from \(\operatorname{periodiccoupledimportancestry}(T, P, m, N, g)\).
\end{proof}
\begin{theorem}
\label{ll:observations:seed-coupled-import-at-arrival}
\(\operatorname{leadingadvance}(14, 5, 3)\) and \(\operatorname{centercleanthrough}(14, 5, 3) \land \operatorname{whitecentergap}(5, 2) \land \operatorname{center}(5) = \operatorname{true} \land \operatorname{rightfront}(14, 6, 2) \land \operatorname{headimportpath}(14, 6, 7)\).
\end{theorem}
\begin{proof}
The goal follows from \(\operatorname{seedcoupledimportatarrival}\).
\end{proof}
\begin{theorem}
\label{ll:observations:ray-bit-balance}
For every natural number \(t\) and natural number \(k\) and natural number \(n\), \(\operatorname{bitvalue}(\operatorname{raybit}(t, k, n)) = \operatorname{bitvalue}(\operatorname{raybit}(t, k, 0)) + \operatorname{rayflux}(t, k, n)\).
\end{theorem}
\begin{proof}
The goal follows from \(\operatorname{raybitbalance}(t, k, n)\).
\end{proof}
\begin{theorem}
\label{ll:observations:diagonal-ray-balance}
For every natural number \(Q\) and natural number \(d\) and natural number \(t\) and natural number \(p\), if \(d + p \le t\), then \(\operatorname{leftphase}(Q + Q, t + (t - d - p), p) = \operatorname{diagonalphase}(Q, d, t) + \operatorname{diagonalflux}(Q, d, t, p)\).
\end{theorem}
\begin{proof}
The goal follows from \(\operatorname{diagonalraybalance}(Q, d, t, p)\).
\end{proof}
\begin{theorem}
\label{ll:observations:diagonal-flux-lag-add}
For every natural number \(Q\) and natural number \(R\) and natural number \(d\) and natural number \(t\) and natural number \(p\), \(\operatorname{diagonalflux}(Q + R, d, t, p) = \operatorname{diagonalflux}(Q, d, t, p) + \operatorname{diagonalflux}(R, d, t + Q, p)\).
\end{theorem}
\begin{proof}
The goal follows from \(\operatorname{diagonalfluxlagadd}(Q, R, d, t, p)\).
\end{proof}
\begin{theorem}
\label{ll:observations:diagonal-phase-doubles}
For every natural number \(Q\) and natural number \(z\) and natural number \(t\), if \(\operatorname{zerodiagonal}(z)\), then if \(\operatorname{doublesat}(Q, z)\), then if \(\operatorname{settledby}(z + 1, t)\), then \(\operatorname{diagonalphase}(Q, z + 1, t) = -1 \lor \operatorname{diagonalphase}(Q, z + 1, t) = 1\) and \(\operatorname{diagonalphase}(Q, z + 1, t + Q) = -\operatorname{diagonalphase}(Q, z + 1, t) \land \operatorname{diagonalphase}(Q + Q, z + 1, t) = 0\).
\end{theorem}
\begin{proof}
The goal follows from \(\operatorname{diagonalphasedoubles}(Q, z, t)\).
\end{proof}
\begin{theorem}
\label{ll:observations:diagonal-reset-phase}
For every natural number \(Q\) and natural number \(z\) and natural number \(t\), if \(\operatorname{zerodiagonal}(z)\), then if \(\operatorname{settledby}(z + 2, t)\), then \(\operatorname{diagonalphase}(Q, z + 3, t + 1) = -\operatorname{diagonalphase}(Q, z + 1, t)\).
\end{theorem}
\begin{proof}
The goal follows from \(\operatorname{diagonalresetphase}(Q, z, t)\).
\end{proof}
\begin{theorem}
\label{ll:observations:diagonal-reset-flux-jump}
For every natural number \(Q\) and natural number \(z\) and natural number \(t\), if \(\operatorname{zerodiagonal}(z)\), then if \(\operatorname{settledby}(z + 2, t)\), then if \(z + 3 \le t\), then \(\operatorname{diagonalflux}(Q, z + 3, t + 1, 1) = \operatorname{diagonalflux}(Q, z + 1, t, 1) + 2 \cdot \operatorname{diagonalphase}(Q, z + 1, t)\).
\end{theorem}
\begin{proof}
The goal follows from \(\operatorname{diagonalresetfluxjump}(Q, z, t)\).
\end{proof}
\begin{theorem}
\label{ll:observations:periodic-diagonal-white-head-balance}
For every natural number \(T\) and natural number \(P\) and natural number \(m\) and natural number \(d\) and natural number \(t\), if \(\operatorname{centerperiodic}(T, P)\), then if \(d + 1 \le t\), then if \(T \le t + (t - d - 1)\), then if \(\operatorname{center}(t + (t - d - 1)) = \operatorname{false}\), then \(\operatorname{diagonalphase}(m \cdot P, d, t) + \operatorname{diagonalflux}(m \cdot P, d, t, 1) = (1 - 2 \cdot \operatorname{bitvalue}(\operatorname{center}(t + (t - d - 1) + 1))) \cdot \operatorname{rightphase}(m \cdot P + m \cdot P, t + (t - d - 1), 1)\).
\end{theorem}
\begin{proof}
The goal follows from \(\operatorname{periodicdiagonalwhiteheadbalance}(T, P, m, d, t)\).
\end{proof}
\begin{theorem}
\label{ll:observations:periodic-diagonal-import-balance}
For every natural number \(T\) and natural number \(P\) and natural number \(m\) and natural number \(d\) and natural number \(t\) and natural number \(start\), if \(\operatorname{centerperiodic}(T, P)\), then if \(d + 1 \le t\), then if \(T \le start\), then if \(\operatorname{headimportpath}(m \cdot P + m \cdot P, start, t + (t - d - 1))\), then if \(\operatorname{center}(t + (t - d - 1)) = \operatorname{false}\), then \(\operatorname{diagonalphase}(m \cdot P, d, t) + \operatorname{diagonalflux}(m \cdot P, d, t, 1) = (1 - 2 \cdot \operatorname{bitvalue}(\operatorname{center}(t + (t - d - 1) + 1))) \cdot (\operatorname{centersign}(start, t + (t - d - 1) - start) \cdot \operatorname{rightphase}(m \cdot P + m \cdot P, start, 2))\).
\end{theorem}
\begin{proof}
The goal follows from \(\operatorname{periodicdiagonalimportbalance}(T, P, m, d, t, start)\).
\end{proof}
\begin{theorem}
\label{ll:observations:coupled-right-front-exact}
For every natural number \(Q\) and natural number \(b\) and natural number \(k\) and natural number \(start\) and natural number \(q\), if \(0 < Q\), then if \(3 \le k\), then if \(\operatorname{leadingadvance}(Q, b, k)\), then if \(\operatorname{centercleanthrough}(Q, b, k)\), then if \(\operatorname{whitecentergap}(b, k - 1)\), then if \(\operatorname{headimportpath}(Q, start, b + (k - 1))\), then if \(1 \le q\), then if \(q + 1 < k\), then \(\operatorname{rightfront}(Q, b + 1, k - 1)\) and \(\operatorname{rightcell}(b + 1, q) = \operatorname{false} \land \operatorname{rightcell}(b + 1 + Q, q) = \operatorname{false}\).
\end{theorem}
\begin{proof}
The goal follows from \(\operatorname{coupledrightfrontexact}(Q, b, k, start, q)\).
\end{proof}
\begin{theorem}
\label{ll:observations:coupled-forced-white-staircase}
For every natural number \(Q\) and natural number \(b\) and natural number \(k\) and natural number \(start\) and natural number \(i\), if \(0 < Q\), then if \(3 \le k\), then if \(\operatorname{leadingadvance}(Q, b, k)\), then if \(\operatorname{centercleanthrough}(Q, b, k)\), then if \(\operatorname{whitecentergap}(b, k - 1)\), then if \(\operatorname{headimportpath}(Q, start, b + (k - 1))\), then if \(1 \le i\), then if \(i + 1 < k\), then \(\operatorname{rightfront}(Q, b + i, k - i)\) and \(\operatorname{rightcell}(b + i, k - i - 1) = \operatorname{false} \land \operatorname{rightcell}(b + i + Q, k - i - 1) = \operatorname{false}\).
\end{theorem}
\begin{proof}
The goal follows from \(\operatorname{coupledforcedwhitestaircase}(Q, b, k, start, i)\).
\end{proof}
\begin{theorem}
\label{ll:observations:depth-two-cancellation-cases}
For every natural number \(Q\) and natural number \(b\) and natural number \(start\), if \(\operatorname{leadingadvance}(Q, b, 2)\), then if \(\operatorname{centercleanthrough}(Q, b, 2)\), then if \(\operatorname{headimportpath}(Q, start, b + 1)\), then \(\operatorname{center}(b) = \operatorname{true}\) and \(\operatorname{center}(b + 1) = \operatorname{false} \land (start < b \land \operatorname{rightdefect}(Q, b, 1) = \operatorname{true} \land \operatorname{rightimportdifference}(Q, b) = \operatorname{false} \land \operatorname{leftphase}(Q, b, 2) = (1 - 2 \cdot \operatorname{bitvalue}(\operatorname{center}(b + 2))) \cdot \operatorname{rightphase}(Q, b, 1) \lor start = b \land \operatorname{rightdefect}(Q, b, 1) = \operatorname{false} \land \operatorname{rightcell}(b, 1) = \operatorname{false} \land \operatorname{rightcell}(b + Q, 1) = \operatorname{false} \land \operatorname{rightdefect}(Q, b, 2) = \operatorname{true} \land \operatorname{rightimportdifference}(Q, b) = \operatorname{true} \land \operatorname{leftphase}(Q, b, 2) = (1 - 2 \cdot \operatorname{bitvalue}(\operatorname{center}(b + 2))) \cdot \operatorname{rightphase}(Q, b, 2))\).
\end{theorem}
\begin{proof}
The goal follows from \(\operatorname{depthtwocancellationcases}(Q, b, start)\).
\end{proof}
\begin{theorem}
\label{ll:observations:coupled-import-birth-or-depth-two-survival}
For every natural number \(Q\) and natural number \(b\) and natural number \(k\) and natural number \(start\), if \(\operatorname{leadingadvance}(Q, b, k)\), then if \(\operatorname{centercleanthrough}(Q, b, k)\), then if \(\operatorname{whitecentergap}(b, k - 1)\), then if \(\operatorname{headimportpath}(Q, start, b + (k - 1))\), then \(start = b + (k - 2)\) or \(k = 2\) and \(start < b \land \operatorname{rightdefect}(Q, b, 1) = \operatorname{true}\).
\end{theorem}
\begin{proof}
The goal follows from \(\operatorname{coupledimportbirthordepthtwosurvival}(Q, b, k, start)\).
\end{proof}
\begin{theorem}
\label{ll:observations:coupled-black-root-prefix}
For every natural number \(Q\) and natural number \(b\) and natural number \(k\) and natural number \(start\) and natural number \(q\), if \(0 < Q\), then if \(3 \le k\), then if \(\operatorname{leadingadvance}(Q, b, k)\), then if \(\operatorname{centercleanthrough}(Q, b, k)\), then if \(\operatorname{center}(b) = \operatorname{true}\), then if \(\operatorname{whitecentergap}(b, k - 1)\), then if \(\operatorname{headimportpath}(Q, start, b + (k - 1))\), then if \(q + 2 < k\), then \(\operatorname{rightcell}(b, q) = \operatorname{true}\) and \(\operatorname{rightcell}(b + Q, q) = \operatorname{true}\).
\end{theorem}
\begin{proof}
The goal follows from \(\operatorname{coupledblackrootprefix}(Q, b, k, start, q)\).
\end{proof}
\begin{theorem}
\label{ll:observations:seed-depth-two-survival}
\(\operatorname{leadingadvance}(4, 5, 2)\) and \(\operatorname{centercleanthrough}(4, 5, 2) \land \operatorname{headimportpath}(4, 2, 6) \land \operatorname{rightdefect}(4, 5, 1) = \operatorname{true} \land \operatorname{rightimportdifference}(4, 5) = \operatorname{false}\).
\end{theorem}
\begin{proof}
The goal follows from \(\operatorname{seeddepthtwosurvival}\).
\end{proof}
\end{document}
